\documentclass{birkjour}
\usepackage{graphicx}
\usepackage{amsmath}
\usepackage{amssymb}
\usepackage{url}

% extended math fonts
\usepackage{yhmath}

\usepackage{srctex} % MikTex,Inverse Search in tex files from Yap screen

\newcommand{\ii}{\mathrm{i}} %imaginary unit
\newcommand{\ee}{\mathrm{e}} %exponent
\newcommand{\dd}{\mathrm{d}} %differential
\newcommand{\cl}[2]{\ensuremath{\mathit{Cl}_{#1,#2}}} % clifford algebra notation


%\DeclareFontFamily{U}{mathx}{\hyphenchar\font45}
%\DeclareFontShape{U}{mathx}{m}{n}{
%      <5> <6> <7> <8> <9> <10>
%      <10.95> <12> <14.4> <17.28> <20.74> <24.88>
%      mathx10
%      }{}
%\DeclareSymbolFont{mathx}{U}{mathx}{m}{n}
%\DeclareFontSubstitution{U}{mathx}{m}{n}
%\DeclareMathAccent{\widecheck}{0}{mathx}{"71}
%\DeclareMathAccent{\wideparen}{0}{mathx}{"75}


\newcommand{\reverse}[1]{\widetilde{#1}}
\newcommand{\gradeinverse}[1]{\wideparen{#1}}
\newcommand{\cliffordconjugate}[1]{\wideparen{\widetilde{#1}}}

\def\d{\!\cdot\!} %narrow-space dot product
\def\w{\!\wedge\!} %narrow-space wedge product
\def\D{\cdot} %wide-space dot product
\def\W{\wedge} %wide-space wedge product


\newcommand{\e}[1]{\mathbf{e}_{#1}} % Clifford basis vectors in Cl(3,1)
\newcommand{\Ie}[1]{I\negthinspace\mathbf{e}_{#1}}%bivector in Cl(3,1)
\newcommand{\inv}[1]{\overline{#1}} %space inversion
\newcommand{\ba}{\mathbf{a}} %bold a
\newcommand{\bk}{\mathbf{k}}
\newcommand{\bx}{\mathbf{x}}
\newcommand{\bbf}{\mathbf{f}} %bold f

\def\m#1{\mathsf{#1}} % general multivector

\def\sR{\mathsf{R}} %sans serif Rotation
\def\sH{\mathsf{H}} %sans serif Hamiltonian
\def\sF{\mathsf{F}} %sans serif F
\def\sG{\mathsf{G}} %sans serif G
\def\cB{\mathcal{B}} %bivector


\begin{document}

%Insert here the title, affiliations and abstract:

\title[Inverse of  multivector: Beyond p+q=5]
 {Inverse of  multivector: Beyond p+q=5\\ threshold}

%----------Author 1
\author{A.~Acus}
\address{Institute of Theoretical Physics and Astronomy,\br
Vilnius University,\br Saul{\.e}tekio 3, LT-10257 Vilnius,
Lithuania} \email{arturas.acus@tfai.vu.lt}

%\thanks{This work was completed with the support of our \TeX-pert.}
%----------Author 2
\author{A.~Dargys}

\address{%
Center for Physical Sciences and Technology,\br Semiconductor
Physics Institute,\br Saul{\.e}tekio 3, LT-10257 Vilnius,
Lithuania}

\email{adolfas.dargys@ftmc.lt}
%----------classification, keywords, date
%----------classification, keywords, date
\subjclass{Primary 15A18; Secondary 15A66}

\keywords{{C}lifford algebra, geometric algebra, inverse
multivector, computer-aided theory}

\date{June, 2017}
%----------additions
%\dedicatory{To my boss}


\begin{abstract}
The algorithm of finding inverse multivector (MV) in a symbolic
form is of paramount importance in computational Clifford
geometric algebra (GA). First attempts of inversion of general MV 
(J.P.~Fletcher, 2002) were based on the decomposition of MV into
basis elements. The complexity of such calculations, nonetheless,
grows exponentially with the dimension $n=p+q$ of $Cl_{p,q}$
algebra. The breakthrough occurred 10 years later (P.~Dadbeh,
2011), after grade-negation operation was introduced. This allowed
to write down  explicit and compact inverse MV as a product of
initial MV and its carefully chosen grade-negation counterparts
for all GAs up to dimension $n\le5$.  It this report we show that
the grade-negation self-product method can be extended beyond
p+q=5 threshold if, in addition, properly constructed multilinear
combinations of such MV products are used. In particular, we write
down explicit compact MV inverse formulas for $p+q=6$ algebras.
\end{abstract}

\maketitle


%%%%%%%%%%%%%%%%%%%%%%%%%%%%%%%%
\section{Introduction\label{sec:1}}

The knowledge of how to find the inverse multivector (MV) in
Clifford algebra in a symbolic and coordinate-free form is very
important both from computational and theoretical point of views.
A universal formula for inverse MV allows to write down a fast and
general algorithm for all occasions rather then to resort to
specific symbolic or numerical cases. The inverse of MV can be
used to find explicit solutions of algebraic GA equations with all
the ensuing consequences and applications. A closely related
problem is the normalization of spinors~\cite{Vaz2016,Lundholm09}.
Invertibility also serves as an important criterion on deciding
whether homogeneous versors are the blades~\cite{Bouma2001} which
are essential in numerous geometric constructions.

First attempts of inversion of some specific forms of MVs can be
traced to papers~\cite{Semenov1991, Semenov1993A,Semenov1993B}.
However it was not until 2002 when
J.P.~Fletcher~\cite{Fletcher2002} has suggested some general MV
inverse formulas for low dimensional algebras. The algorithm of
these inverse MVs was based on decomposition of MV into basis
elements, and, as a consequence, resulted in exponential growth of
the complexity of problem along with  Clifford algebra $Cl_{p,q}$
dimension $n=p+q$.

The breakthrough occurred 10 years later when
P.~Dadbeh~\cite{Dadbeh2011} introduced grade-negation operation
that has  allowed  to write down compact and explicit formulas for
MV inverse in a coordinate-free form. The heart of the algorithm
is a geometric product of initial MV and its carefully chosen
grade-negated counterparts that after few iterations eventually
yield a scalar. In fact, the product may be related with a
determinant of matrix representation of general MV, from which the
inverse multivector can be easily extracted  by simply removing
the initial MV, which is always positioned either in the left-most
or right-most side of the product. The described method allows to
find explicit inverse formulas for a general MV in $n$-dimensional
Clifford space  up to dimension $n=p+q\le 5$, independent of GA
signature $(p,q)$. It was rewritten by number of authors using
different notations and forms (see, for
example,~\cite{Shirokov2012,Hitzer2016}). Similar formulas for the
norms of MVs for a few specific Clifford algebras when  $n\le 5$
was also given  few years before~\cite{Lundholm09}. Analysis of
general structure of such formulas was presented
in~\cite{Suzuki2016}. However, until now any attempts to
step across the threshold  $p+q=5$ were unsuccessful.

It this report we show that the grade-negation method can be
extended beyond $p+q=5$ threshold if, in addition, properly
constructed  multilinear combinations of grade-negated MVs are
used. In particular,  we write down explicit MV inverse formulas
for $p+q=6$ algebras having all possible   $(p,q)$ signatures. In
Sec.~\ref{sec:2} the grade-negation method is shortly reviewed and
required  notation is introduced. In Sec.~\ref{sec:3}  the inverse
even MVs that follow from higher grade inverse MVs are considered.
Finally, in Sec.~\ref{sec:4} a general algorithm for construction
 of inverse MV at $p+q=6$ is explained and the obtained coordinate-free formulas
are presented.


%%%%%%%%%%%%%%%%%%%%%%%%%%%%%%%%%%%%%%%%%%%%%%%%%%%%%%%
\section{Grade-negated self-product and inverse of general MV of algebras of dimension $n=p+q\le5$ \label{sec:2}}

Following~\cite{Dadbeh2011} we first introduce a grade-negated
self-product, which  is defined via \textit{grade-$r$ negation
operation}.  Applied to the multivector $\m{A}$ this operation
does what it says, i.~e., it changes the sign of grade-$r$ part in the
MV~$\m{A}$.  Such grade-$r$ negated MV  will be denoted as
$\m{A}_{\bar r}$, with the bar over index designating which of the
grades have opposite signs.  Formally the grade-$r$ negated MV can
be expressed as $\m{A}_{\bar r}=\m{A}-2\langle \m{A}\rangle_r$,
where $\langle \m{A}\rangle_r$ denotes grade-$r$ projection of
multivector  $\m{A}$. In particular, we have
$(\langle\m{A}\rangle_r)_{\bar r}=-\langle\m{A}\rangle_r$.

The following properties are evident  from definition of
grade-negation operation: $(\m{A}_{\bar p})_{\bar r}=\m{A}_{{\bar
p},{\bar r}}=\m{A}_{{\bar r},{\bar p}}$, $\m{A}_{{\bar r},{\bar
r}}=\m{A}$, $(\m{A}+\m{B})_{\bar r}=\m{A}_{\bar r}+\m{B}_{\bar
r}$. However, $(\m{AB})_{\bar r}\ne \m{A}_{\bar r}\m{B}_{\bar r}$.
If $\m{A}$ does not contain grade-$r$ elements then negation
returns the same MV, $\m{A}_{\bar r}=\m{A}$. Commutator with
grade-negated MV then can be expressed as $\m{A} \m{A}_{\bar
r}-\m{A}_{\bar r}\m{A}=2(\langle \m{A}\rangle_r\m{A}-\m{A}\langle
\m{A}\rangle_r)$. The standard reversion denoted as
$\reverse{\m{A}}$, grade inversion $\gradeinverse{\m{A}}$, and
Clifford conjugate
$\cliffordconjugate{\m{A}}=\reverse{\gradeinverse{\m{A}}}$
operations in term of negation, respectively, can be written as
$\reverse{\m{A}}=\m{A}_{{\bar 2},{\bar 3},{\bar 6},{\bar 7},{\bar
{10}},{\bar {11}}\dots}$, $\gradeinverse{\m{A}}=\m{A}_{{\bar
1},{\bar 3},{\bar 5},{\bar 7},{\bar {9}},{\bar {11}}\dots}$ and
$\cliffordconjugate{\m{A}}=\m{A}_{{\bar 1},{\bar 2},{\bar 5},{\bar
6},{\bar {9}},{\bar {10}}\dots}$. For a concrete algebra the
series of indices can be chopped off when the value of the
negation index becomes larger than pseudoscalar index. It should
be noted that grades $0,4,8,12,\dots$ are absent in the above list
of standard GA involutions. From this follows that the inverse of
a general MV cannot be expressed using just standard involutions,
or their combinations. As we shall see this observation also is directly related to the MV inversion problem for $p+q\ge 6$.

By grade-negated self-product we identify the geometric product of
general multivector $\m{A}$ with any number of its grade-negated
counterparts, for example, $\m{A} \m{A}_{\bar r} \m{A}_{\bar
k,\bar l,\bar t} \m{A}_{\bar s, \bar t}\cdots$ or $\m{A}_{\bar r}
\m{A}_{\bar k,\bar l,\bar t} \m{A}_{\bar s, \bar t}\cdots\m{A} $.
Those negated self-products where the initial MV $\m{A}$ stands in
the left-most or in the right-most position will be of special
importance in finding the inverse multivector $\m{A}^{-1}$. It is
obvious, that if we have  some identity written in Clifford
algebra of dimension $n=p+q$ which includes only operations of
geometric product and general grade negation operation  then this
identity will be valid for any $n$-dimensional algebra having
arbitrary signature $(p,q)$. This happens because grade negation
(which is, of course, one of GA involutions) completely
ignores the signature of the algebra.

The explicit inverse formulas which we will construct rely on our
ability to get a scalar using just grade-negated self-products
(or, as we shall see, some multilinear combinations of them for
$n>5$), where initial multivector is a mandatory and located
either in the left-most or right-most position. We shall call in
this way obtained scalar
\begin{equation}\label{discriminant}
s_m=\overbrace{\m{A}\m{A}_{\bar r,\bar t,\ldots}\m{A}_{\bar
s,\ldots}\cdots}^m \end{equation}  \textit{a discriminant of the
MV} if its symbolic expression coincides with the determinant of
the matrix representation of the multivector $\m{A}$. If we succeed in
constructing such a scalar then it is easy to see that left or
right inverse of MV can be obtained simply by removing left or
right  most positioned initial MV $\m{A}$ from the scalar
$s_m=\m{A}\m{A}_{\bar r,\bar t,\ldots}\m{A}_{\bar s,\ldots}\cdots$
and then dividing the result by $s_m$, i.~e., the left inverse of MV $\m{A}$
is
\begin{equation}\label{inverse}
\m{A}^{-1}=\frac{\m{A}_{\bar r,\bar t,\ldots}\m{A}_{\bar
s,\ldots}\cdots}{s_m}.\end{equation}

For low dimensional algebras the construction of the scalar $s_m$
is easy and can be done by a simple rule of thumb: in each of
intermediate steps of construction,  the self-product of two terms
should not repeat previous grades  and at the same time should
eliminate a sufficient number of remaining grades, so that the
midway MV will consist of smaller number of grades. However, the
rule generally does not guarantee that this always will lead to the desired result
(scalar). We will demonstrate that rule works for algebras with
$n\le5$ and fails starting from $n\ge6$. To facilitate the guess
which of grades must  be negated in the self-products we have
written \textit{Mathematica} symbolic
program~\cite{AcusDargys2017} for fast calculation of geometric
products. Examples of such products can be seen in
table~\ref{invall}.

{\bf Two dimensional algebras.} Since one dimensional case is
trivial we shall start from two dimensional algebras $Cl_{2,0}$,
$Cl_{1,1}$ and $Cl_{0,2}$. As mentioned,  the  elimination of
grades  by self-products does not depend on algebra signature,
therefore, we may choose any of algebras. For $n=2$ let it be
$Cl_{2,0}$. Writing $\m{A} = \sum_{J=0}^{2^n-1} a_J\e{J}$, where
the multi-index $J$ covers all orthonormal base elements arranged
in inverse degree lexicographic order, i.~e., the lowest grade
elements appear first while elements of the same grade are ordered
lexicographically. In the considered example the multi-index takes
the following explicit values $J=[\{\},\{1\},\{2\},\{1,2\}]$. We
can check that self-product $\m{A} \m{A}_{\bar 1,\bar 2}$
immediately gives the required scalar $s_2=a_{\{\}}^2 - a_{1}^2 -
a_{2}^2 + a_{1,2}^2$. In the standard notation it is more
convenient to rewrite the coefficient indices as $s_2=a_{0}^2 -
a_{1}^2 - a_{2}^2 + a_{3}^2$, which means that  in front of base
scalar~$1$ the index in  coefficient  $a_{0}$ is zero, the
coefficients at vectors run from $1$ to $n$ while numerical
indices of higher grades increase monotonically up to $2^n$. In
general, when programming, for a coefficient in front of grade-$r$
base it is convenient to start the element index
enumeration by $\sum_{k=0}^{r-1} (\begin{smallmatrix}n\\
k\end{smallmatrix})$ and to end by
$\Bigl(\sum_{k=0}^{r-1} (\begin{smallmatrix}n\\
k\end{smallmatrix})\Bigr)+(\begin{smallmatrix}n\\
r\end{smallmatrix})-1$.  Here $(\begin{smallmatrix}n\\
k\end{smallmatrix})$ denotes the binomial coefficient. For
example, for \cl{3}{0}, which has three vectors and three
bivectors, these expressions give $1$ and $3$ for first/last
vector indices, and $4$ and $6$ for first/last bivector indices.
Taking into account that binomial with negative index vanishes,
this convention enumerates all coefficients of MV and allows an
easy transition to more standard notation in a consistent way, for
example, $\m{A}=a_0 + a_1\e{1} + a_2\e{2} + a_3\e{12}$ for initial
MV and $\m{A}_{\bar 1,\bar 2}=a_0 - a_1\e{1}- a_2\e{2}- a_3\e{12}$
for grade-negated MV.  This notation  will be used throughout the
paper. Note, that indices of base elements will never be renamed
and always be referred to as multi-indices. If no ambiguity will
arise, we also shall delete the separating commas between indices
in base elements. Before  going to higher dimensional algebras few
comments are in place.

 (1) The explicit form of discriminant $s_m$ comprises all coefficients $a_i$ of multivector~$\m{A}$.

 (2) The condition for a MV to have inverse is determined by
 discriminant which must be unequal to zero, $s_m\ne 0$. This imposes
 the requirement between MV coefficients.

(3) Because reverse leaves the scalar (discriminant) invariant,
the above formula for $n=2$ can be rewritten as
$s_2=\m{A} \m{A}_{\bar 1,\bar 2}=\m{A}_{\bar 1,\bar 2}\m{A}$, from
which the left/right inverse formulas follow,
$\m{A}^{-1}=\m{A}_{\bar 1,\bar 2}/s_2= \m{A}_{\bar 1,\bar
2}/(\m{A}\m{A}_{\bar 1,\bar 2})=\m{A}_{\bar 1,\bar 2}/(\m{A}_{\bar
1,\bar 2}\m{A})$. Since  $\m{A}_{\bar 1,\bar
2}=\cliffordconjugate{\m{A}}$, for \cl{2}{0} algebra the
negated MV can be expressed through standard involutions.

 (4) The discriminant expression for $n=2$ contains only two MVs.
 From 8-periodicity table~\cite{Lounesto97} follows that GA
 algebras \cl{2}{0}, \cl{1}{1} and \cl{0}{2}
 can be represented  by  $2\times2$ real matrices $\mathbb{R}(2)$ or quaternion~$\mathbb{H}$.
 Thus we have that both the determinant of real matrix and the product of two MVs are
 quadratic in MV coefficients $a_i$. Below we shall see that in all cases the highest degree  of coefficients in the discriminant
 always matches the degree of determinant of real matrix representation of $\m{A}$.
 This suggests that the number of MVs needed in the grade-negated self-product can be determined
 from the size of real irreducible matrix representation in 8-periodicity table.


 (5) The eliminated grades are grade-negated in one of MVs in the self-product.
 In the considered  example the grades $1$ and $2$ have been eliminated simultaneously
 in  $\m{A}_{\bar 1,\bar 2}$. As we shall see this rule is valid for discriminants
 up to $n\le5$.
 For $n\ge6$ exceptions from the rule appear due to multiples of grade-$4$ and
 because of reappearance of some previously eliminated grades.

{\bf Three dimensional algebras.}  Our goal is to eliminate as
many grades as possible by forming suitable self-products. What is
important is that the remaining grades in the self-products should
not regenerate already eliminated grades. As another example, let
us  calculate the inverse of $Cl_{2,1}$ algebra. It is easy to
check that the product of $\m{A}=a_{0}+a_{1} \e{1}+a_{2}
\e{2}+a_{3} \e{3}+a_{4} \e{12}+a_{5} \e{13}+a_{6} \e{23}+a_{7}
\e{123}$ with $\m{A}_{\bar 1,\bar 2}$ lacks grades $1$ and $2$.
The result is the new multivector $\m{B}=\m{A} \m{A}_{\bar 1,\bar
2}= b_{0}+b_{7} \e{123}$ which  consists of scalar and grade-$3$
element. In particular, $b_0=a_{0}^2 - a_{1}^2 - a_{2}^2 + a_{3}^2
+ a_{4}^2 - a_{5}^2 - a_{6}^2 + a_{7}^2$ and $b_7=-2 a_{3} a_{4} +
2 a_{2} a_{5} - 2 a_{1} a_{6} + 2 a_{0} a_{7}$. Repeating  grade
negation procedure one finds that grade-$3$ part can be removed
too, and we obtain the left-handed discriminant $s_4=\m{B}
\m{B}_{\bar 3}=\m{A} \m{A}_{\bar 1,\bar 2}(\m{A} \m{A}_{\bar
1,\bar 2})_{\bar 3}=b_0^2 - b_7^2$. Reversion of $s_4$ then
immediately yields the $\m{A}$-right-handed discriminant
$s_4^\prime=(\m{A}_{\bar 1,\bar 2}\m{A})_{\bar 3} \m{A}_{\bar
1,\bar 2}\m{A}$ that contains four MVs. Of course, for different
3D algebras the same formula will give different fully expanded
expressions, which will differ in signs of individual
coefficients~$a_i$. Note, that now the discriminant is determined
by two terms, $b_0^2$ and $b_7^2$, the difference of which should
not be zero for invertible multivector $\m{A}$ to exist. The
condition $b_0^2-b_7^2\ne 0$ exactly matches the MV invertibility
condition, which  according to 8-periodicity table may be obtained
by calculating the determinant of matrix representation of MV,
namely, the determinant of $4\times 4$ real matrices
${^2\mathbb{R}(2)}$.

{\bf Four dimensional algebras.} In $n=4$ case, trying to
eliminate as much grades as possible  we can proceed in two
alternative ways. Firstly, we can negate simultaneously grades $1$
and $2$ and then in the next step the grades $3$ and $4$.
Alternatively  we can eliminate $2$ and $3$, and then $1$ and $4$
grades. Both choices are valid. However, if we choose  $1$ and
$3$, and  $2$ and $4$ grade combinations, neither one will  do the
job. Thus, we find the following formulas for the left-hand
discriminant: $s_4= \m{A}\m{A}_{\bar{1},\bar{2}}(\m{A}
\m{A}_{\bar{1},\bar{2}})_{\bar{3},\bar{4}}= \m{A}
\m{A}_{\bar{2},\bar{3}}
(\m{A}\m{A}_{\bar{2},\bar{3}})_{\bar{1},\bar{4}}$. The right-hand
version  is similar,
$s_4^\prime=(\m{A}_{\bar{2},\bar{3}}\m{A})_{\bar{1},\bar{4}}\m{A}_{\bar{2},\bar{3}}
\m{A}=
(\m{A}_{\bar{1},\bar{2}}\m{A})_{\bar{3},\bar{4}}\,\m{A}_{\bar{1},\bar{2}}
\m{A}$.  The appearance of symbol $\m{A}$ as much as four times in
the products $s_4$ and  $s_4^{\prime}$ is again what we expect
from the matrix representation for $n=4$ algebras, namely, the
highest degree of determinant polynomial should be of degree $4$
as well, in accord with $4\times 4$ matrix representation
$\mathbb{R}(4)$. Despite different forms, expanded explicit
expressions for $s_m$ and $s_m^\prime$ in fact always are equal,
$s_m=s_m^\prime$.


{\bf Five dimensional algebras.} This is the largest dimension
$n=p+q=5$ at which consecutive grade elimination algorithm, that
factorizes the discriminant into product of initial and  negated
MV's, works. The grade elimination sequence is similar: (1)
Eliminate simultaneously grades $2$ and $3$; (2) Then eliminate
grades $1$ and $4$; (3) Finally eliminate grade $5$. Apart from
the last step this sequence is exactly the same as the second
alternative of $n=4$ case. The final result is
$s_8^\prime=(\m{F}\m{A})_{\bar{5}}\m{F}\m{A}$ with
$\m{F}=(\m{A}_{\bar{2},\bar{3}}\m{A})_{\bar{1},\bar{4}}
\m{A}_{\bar{2},\bar{3}}$ for right-hand sided form of the
discriminant and $s_8=\m{A}\m{D}\, (\m{A}\m{D})_{\bar{5}}$ with
$\m{D}=\m{A}_{\bar{2},\bar{3}}
(\m{A}\m{A}_{\bar{2},\bar{3}})_{\bar{1},\bar{4}}$ for left-hand
discriminant. Matrix representation of five dimensional Clifford
algebras corresponds to   real ${^2\mathbb{R}(4)}$ matrix the the
determinant of which is a polynomial of degree $8$. This is in
accord with the structure of discriminant~$s_8=s_8^\prime$ where
the $\m{A}$ occurs as much as eight times.


Can this grade elimination procedure be extended beyond $n=5$? The
short answer is 'yes' if, as shown below, we allow multilinear
combinations of grade-negated self-products with properly chosen
numerical coefficients. Before describing this case let us derive
some useful formulas for even subalgebras when  $n\le6$. The even
subalgebras are related with the spinor groups that are very
important in the quantum mechanics~\cite{Vaz2016}.

%%%%%%%%%%%%%%%%%%%%%%%%%%%%%%%%%%%%%%%%%%%%%%%%%%%%%%%
\section{Inverse of even multivectors\label{sec:3}}


\begin{table}
\[\begin{array}{ll}
Cl_{p,q}&\qquad  \bigl(\m{A}^{+}\bigr)^{-1} \\[5pt]\hline
p+q=2&\qquad \frac{B \bigl(\m{A}^{+} B\bigr)} {\bigl(\m{A}^{+}
B\bigr)^2}, \qquad B =(\e{i})_{\bar{1}}=-\e{i}
\\[4pt]
  p+q=3 &\qquad \frac{C \bigl(\m{A}^{+} C\bigr)}
{\bigl(\m{A}^{+} C\bigr)^2}, \qquad C = (\m{A}^{+}
\e{i})_{\bar{1}}
\\[4pt]
  p+q=4&\qquad \frac{C \bigl(\m{A}^{+} C\bigr)}
{\bigl(\m{A}^{+} C\bigr)^2}, \qquad  C = (\m{A}^{+}
\e{i})_{\bar{1}}
\\[4pt]
 p+q=5&\qquad \frac{D \bigl(\m{A}^{+} D\bigr)}
{\bigl(\m{A}^{+} D\bigr)^2}, \qquad  D = (\m{A}^{+}
\e{i})_{\bar{3}} \bigl(\m{A}^{+} (\m{A}^{+}
\e{i})_{\bar{3}}\bigr)_{\bar{1}}
\\[4pt]
p+q=6&\qquad \frac{D \bigl(\m{A}^{+} D\bigr)_{\bar{6}}}
{\bigl(\m{A}^{+} D\bigr)\bigl(\m{A}^{+} D\bigr)_{\bar{6}}}, \quad
D = (\m{A}^{+} \e{i})_{\bar{3}} \bigl(\m{A}^{+} (\m{A}^{+}
\e{i})_{\bar{3}}\bigr)_{\bar{1}}
\end{array}
\]
\caption{\label{inveven} Explicit formulas of inverse of even
multivectors in a coordinate-free form for Clifford algebras of
dimension $p+q\le 6$. The quantities $\bigl(\m{A}^{+}
B\bigr)^2/\e{i}^2$, $\bigl(\m{A}^{+} C\bigr)^2/\e{i}^2$,
$\bigl(\m{A}^{+} D\bigr)^2/\e{i}^4$ and ${\bigl(\m{A}^{+}
D\bigr)\bigl(\m{A}^{+} D\bigr)_{\bar{6}}}/\e{i}^4$  are the
discriminants of respective multivectors. Here explicit
$\e{i}^2=\pm 1$ and $\e{i}^4=+1$ factors take into account the
correct sign and norm once the orthonormal base vector $\e{i}$ is replaced by arbitrary unnormalized and nonisotropic vector $v$.}
\end{table}


When  MV consists of even grade elements only, i.~e.
$\m{A}^{+}=\langle\m{A}\rangle_{0}+\langle\m{A}\rangle_{2}+\langle\m{A}\rangle_{4}+\cdots$,
simpler formulas for inverse MVs can be constructed as shown in
table~\ref{inveven}. The given formulas in no way do restrict the
general MV case. The base vectors $\e{i}$ in these formulas play
the role of dummy variables and therefore may be replaced by
different basis vectors, or even by any linear combination of base
vectors of nonzero magnitude, i.~e., by any nonisotropic vector.
By reversing numerators and denominators right-hand side formulas
can be straightforwardly obtained.

It is interesting to observe that in contrast to general case at
$p+q=6$ there exists a single self-negated product  for an even
inverse MV (see table~\ref{inveven}) that is not a multilinear
combination. This property is to be expected (compare with  a
general multivector form for $p+q=5$ in table~\ref{invall}),
because there exists the well-known isomorphism between the even
subalgebra of one dimension larger algebra  and the full lower
dimensional Clifford algebra,
\begin{equation}\label{iso}
\cl{p}{q+1}^+\cong\cl{p}{q},\quad\cl{p+1}{q}^+\cong\cl{q}{p}\,.
\end{equation}
Because we can write inverse MV as a single self-negated product
for dimension $p+q=5$ (see table~\ref{invall}), it is not
surprising that according to isomorphisms~\eqref{iso} we can do
this for even subalgebra of dimension $p+q=6$. Unfortunately, as
we shall see in section~\ref{sec:4} this property does not extend
to the full six dimensional Clifford algebra.

%%%%%%%%%%%%%%%%%%%%%%%%%%%%%%%%%%%%%%%%%
\section{Inverse of general MV in $p+q=6$ dimensional Clifford algebra\label{sec:4}}
\subsection{Insufficiency of single self-negated product in $p+q=6$ case\label{sec:4a}}

Let us take $Cl_{6,0}$ algebra. Using \textit{Mathematica} GA
package~\cite{AcusDargys2017}, after some experimentation with a
pair of self-negated product formed from a general MV it is not
difficult to discover that we can eliminate simultaneously either
grades $1$, $2$, $5$ and $6$ (then $0$, $3$ and $4$ grades
survive) or, alternatively,  the grades $2$, $3$ and $6$ (then
grades $0$, $1$, $4$ and $5$ persist). It should be noted that in
both cases the grade $4$ remains and, therefore, cannot be
eliminated by the method of simple self-negated product used till
now. This conclusion strictly follows from an attempt to simultaneously
nullify all coefficients of grade-$4$ base elements using all
$2^6=64$ possible combinations of grade negations in the two term
product. The grade-$4$ therefore is distinct from the other grades and
deserves special examination. It turns out that there exists a
subalgebra with base  formed by elements $\{1, \e{1256},
\e{1346},\e{2345}\}$ such that any self product of multivector
$\m{B}=a_{0}+a_{47} \e{1256}+a_{49} \e{1346}+a_{52} \e{2345}$  by
grade-negated $\m{B}$ will yield new MV having at least one
grade-$4$ base element present. This is the reason why such single
self-product fails in eliminating grade-$4$ part and, as we shall
see, one has to use a combination of at least a pair of
self-products to eliminate such an infinite loop.

\subsection{Multilinear combination of self-products\label{sec:4b}}

Before looking for coefficients of multilinear combination of two
self-negated products, first, we shall note that the determinant
of matrix representation of general MV that consists of scalar and general
$4$-th grade element can be writen as $s_8=(s_4)^2$. So we can always take square
root of it. It means that instead of $m=8$ in
equations~\eqref{discriminant} just a half of them is enough to
write down the discriminant for restricted to scalar + grade-$4$
multivector, i.~e., we can look for the self-negated product which
contains only four independent multivectors (the related
discriminant of multilinear combination of two terms each
consisting of geometric product of $4$ multivectors will be
denoted as $s_{4+4}$). By forming a prototype of such a geometric product of the
form $s_{4+4}=s_{4f} +s_{4g}$ we have
\begin{equation}\begin{split}
 &s_{4f}+s_{4g}=\\
 &b_1 B f_5\Bigl(f_4(B) f_3\bigl(f_2(B) f_1(B)\bigr)\Bigr)
 +b_2 B g_5\Bigl(g_4(B) g_3\bigl(g_2(B)
  g_1(B)\bigr)\Bigr)\label{s1},
\end{split}\end{equation}
where $f_j, g_j$ denotes yet undetermined grade-$4$ negations and
$b_1, b_2$ are unknown scalar coefficients of the multilinear
combination. The above  choice of the discriminant, of course, is
not unique. The other three possible choices (see below), however, yield the same answer.

Once negation search pattern was fixed we can calculate explicit
symbolic form of square root of the determinant of matrix
representation of the multivector $\m{B}=a_{0}+a_{47}
\e{1256}+a_{49} \e{1346}+a_{52} \e{2345}$ and compare it with
 all possible negations of grade-$4$ substituted for $f_j$ and $g_j$ in~\eqref{s1}.
The involutions of grade-$4$ can be modelled as a multiplication
by unknown coefficient $p_{4jk}$, where the index $j$ denotes the
involution number $f_j,g_j$ in the self-product and $k$ is the
term in multilinear combination (in $p+q=6$ case $k=1$ for $f$ and
$k=2$ for $g$). Later, when we shall deal with negations of other
grades the first index $i$ in $p_{ijk}$ will indicate possible
involuted grade-$i$. For the moment the index $i$ is fixed to $4$,
which corresponds  to current grade of $\m{B}$ (we will ignore
negations of scalar, because it is equivalent to negation of all
other MV grades). The coefficients $p_{ijk}$ acquire the values
$\pm 1$ only, where $-1$ means that involution that changes sign
of grade $i$ is applied and $+1$ means that grade-$i$ remains
unchanged.

In the considered case, comparison with the square root of
determinant of matrix representation of $\m{B}$ yields the
following two sets of solutions, $\{b_1=-2/3,\; b_2=-1/3,\;
p_{411}=-1,\; p_{412}=1,\; p_{421}=-1,\; p_{422}=1,\;
p_{431}=-1,\; p_{432}=-1,\; p_{441}=-1,\; p_{442}=1,\;
p_{451}=-1,\; p_{452}=1\}$ and $\{b_1=-1/3,\; b_2=-2/3,\;
p_{411}=1,\; p_{412}=-1,\; p_{421}=1,\; p_{422}=-1,\;
p_{431}=-1,\; p_{432}=-1,\; p_{441}=1,\; p_{442}=-1,\;
p_{451}=1,\; p_{452}=-1\}$. It is enough to take only one of them,
because the other, in fact,  represents the term interchange in
multilinear combination. It is easy to check that the above
solution though was computed from highly simplified MV (scalar
plus three grade-$4$ base elements) works flawlessly for more
general MV (scalar plus any number of grade-$4$ elements) $\m{C}=
a_0+a_{42} \e{1234}+a_{43} \e{1235}+a_{44} \e{1236}+a_{45}
\e{1245}+a_{46} \e{1246}+a_{47} \e{1256}+a_{48} \e{1345}+a_{49}
\e{1346}+a_{50} \e{1356}+a_{51} \e{1456}+a_{52} \e{2345}+a_{53}
\e{2346}+a_{54} \e{2356}+a_{55} \e{2456}+a_{56} \e{3456} $. This
is what one expects, because grade involution acts on all equal
grade elements in exactly identical way. This considerably speeds
up search procedure for involutions, because we always can choose
truncated or shorter  MVs when trying to balance calculation
complexity against speed.

We note that in $s_{4+4}=s_{4f} +s_{4g}$, general form of
negated-product $s_{4f}\sim B f_5\Bigl(f_4(B) f_3\bigl(f_2(B)
f_1(B)\bigr)\Bigr)$ could be replaced by similar products with
interchanged negations, $s_{4f^\prime}\sim B
f_5^\prime\Bigl(f_3^\prime\bigl(f_4^\prime(B)
f_2^\prime(B)\bigl)f_1^\prime(B)\Bigl)$, and similarly for~$g_j$.
We have explicitly verified that in the case of scalar + general
grade-$4$ multivector all four possible variants of discriminant,
i.~e. $(s_{4f}+s_{4g})(s_{4f}+s_{4g})$,
$(s_{4f}+s_{4g^\prime})(s_{4f}+s_{4g^\prime})$,
$(s_{4f^\prime}+s_{4g})(s_{4f^\prime}+s_{4g})$ and
$(s_{4f^\prime}+s_{4g^\prime})(s_{4f^\prime}+s_{4g^\prime})$ yield
the same solution, i.~e.  a single pair (multilinear combination
of two terms) of self-negated products that eliminates grade-$4$
elements from the MV may be selected. We will note that
grade-negation of other grades in equation~\eqref{s1} is not
unique and, in general, provides numerous possibilities. In the
end this will yield different formulas for inverse MV, similarly
as it was for algebras $p+q=4$.


Once the coefficients $b_k$ in equation~\eqref{s1} and arrangement
of involutions $f_j$ and $g_j$ that ensure that the problematic
grade-$4$ elimination were successfully found, we can begin
the search in exactly the same way for concrete involutions that
consecutively will eliminate all other grades. It was already
mentioned that in equations of type~\eqref{s1} either grades $1$,
$2$, $5$ and $6$ (then grades $0$, $3$ and $4$ remain), or,
alternatively grades $2$, $3$ and $6$ (then grades $0$, $1$, $4$
and $5$ remain) can be simultaneously eliminated. All in all there
are $\frac{6!}{2!4!}+\frac{6!}{3!3!} +1=36$ elements of grades
$2$, $3$ and $6$ which can be eliminated simultaneously. This is
more than $28$ total elements of grades $1$, $2$, $5$ and $6$ that need to be eliminated. Therefore, if we replace $\m{B}$ in
equation~\eqref{s1} by self-negated product $\m{A}\m{A}_{{\bar
2},{\bar 3},{\bar 6}}$ we then are left to deal with self-negated
product of four multivectors of grades $0$, $1$, $4$, $5$ only,
from which we can ignore  the grades $0$ and $4$. Thus, repeating
the trick of finding remaining undetermined involutions of grades
$1$, $2$, $5$ we arrive to a number of valid solutions for
coefficients $p_{1jk},p_{2jk}$ and $p_{5jk}$. In order to speed up
the derivation we had used the multivector $ a_{0}+a_{1}
\e{1}+a_{22} \e{123}+a_{47} \e{1256}$ first. The obtained result
then was checked symbolically using the most general $Cl_{6,0}$
multivector. There are $320$ such solutions for the above
considered case, which include different numbers and combinations
of negation operations. Among them there is just one solution,
which applies negation to the minimal number of MVs. It was found
that  minimal number of negated MVs is $10$. However, from
computation point of view the solution with minimal number of
negations isn't optimal. Indeed, while the answer is correct, some
grades cancel out in the last moment, during subtraction between 
terms $s_{4f}$ and $s_{4g}$ in~\eqref{s1} (actually we
had to rename $s_{4f}+s_{4g}\rightarrow s_{8f}+s_{8g}$ now, because $\m{B}$
was replaced by product $\m{A}\m{A}_{{\bar 2},{\bar 3},{\bar 6}}$
and formula~\eqref{s1} now contains multilinear combination of~8 multiplied MVs, which corresponds to polynomial degree of full, now unfactorizable, determinant). Using more involutions we can force these
grades to be eliminated in each of terms separately. Then only
grade-$4$ remains to be cancelled out by the multilinear
combination itself.

After  the discriminant has been found  the inverse multivector
can be immediately written using equation~\eqref{inverse}. The
results for inverse MVs for algebras $p+q\le 6$ are presented in
table~\ref{invall}.
\begin{table}
\[\begin{array}{ll}
\cl{p}{q}&\qquad  \m{A}^{-1} \\[5pt]\hline\\[-5pt]
p+q=0&\qquad\frac{\m{B}}{\m{A}\m{B}},\qquad\qquad\qquad\qquad\m{B}=1\\[4pt]
%
 p+q=1&\qquad
\frac{\m{B}(\m{A}\m{B})_{\bar{1}}}{\m{A}\m{B}\,(\m{A}\m{B})_{\bar{1}}}
= \frac{\gradeinverse{\m{A}}}{\m{A} \gradeinverse{\m{A}}},\quad\qquad\m{B}=1\\[4pt]
%
  p+q=2&\qquad \frac{\m{C}}{\m{A}\m{C}}=\frac{\cliffordconjugate{\m{A}}}{\m{A} \cliffordconjugate{\m{A}}}
,\qquad\qquad\quad\m{C}=\m{A}_{\bar{1},\bar{2}}\\[4pt]
%
  p+q=3 &\qquad
 \frac{\m{C}(\m{A}\m{C})_{\bar{3}}}{\m{A}\m{C}\,(\m{A}\m{C})_{\bar{3}}}
=\frac{\cliffordconjugate{\m{A}}\gradeinverse{\m{A}}\reverse{\m{A}}}{\m{A}
\cliffordconjugate{\m{A}}\gradeinverse{\m{A}}\reverse{\m{A}}}
,\qquad\m{C}=\m{A}_{\bar{1},\bar{2}}\\[4pt]
%
 p+q=4&\qquad\frac{\m{D}}{\m{A}\m{D}},\qquad\qquad\qquad\qquad
\m{D}=\m{A}_{\bar{2},\bar{3}}
(\m{A}\m{A}_{\bar{2},\bar{3}})_{\bar{1},\bar{4}}\quad \\
   &\qquad\qquad\qquad\qquad\qquad\textrm{or}\quad \m{D}=\m{A}_{\bar{1},\bar{2}}(\m{A}
\m{A}_{\bar{1},\bar{2}})_{\bar{3},\bar{4}}  \\[4pt]
%
 p+q=5&\qquad\frac{\m{D}(\m{A}\m{D})_{\bar{5}}}{\m{A}\m{D}\,
(\m{A}\m{D})_{\bar{5}}},\qquad\qquad\qquad
\m{D}=\m{A}_{\bar{2},\bar{3}}
(\m{A}\m{A}_{\bar{2},\bar{3}})_{\bar{1},\bar{4}}\\[4pt]
%
p+q=6&\qquad\frac{\m{G}}{\m{A}\m{G}},\qquad\m{G}=\frac{1}{3} \m{A}_{\bar{2},\bar{3},\bar{6}} \Bigl(\m{H} (\m{H} \m{H})_{\bar{1},\bar{4},\bar{5}} + 2 \bigl(\m{H}_{\bar{4}} (\m{H}_{\bar{4}} \m{H}_{\bar{4}})_{\bar{1},\bar{4},\bar{5}}\bigr)_{\bar{4}}\Bigr)\\
     &\quad\textrm{or}\quad \m{G}=\frac{1}{3} \m{A}_{\bar{2},\bar{3},\bar{6}} \Bigl(\bigl(\m{H} (\m{H} \m{H}_{\bar{1},\bar{5}})_{\bar{4}}\bigr)_{\bar{1},\bar{5}} + 2 \bigl(\m{H}_{\bar{4},\bar{5}} (\m{H}_{\bar{4},\bar{5}} \m{H}_{\bar{1},\bar{4}})_{\bar{4}}\bigr)_{\bar{1},\bar{4}}\Bigr)\\
&\qquad \qquad \qquad\textrm{with}\quad \m{H}=\m{A} \m{A}_{\bar{2},\bar{3},\bar{6}}=\m{A}\reverse{\m{A}}
\end{array}
\]
\caption{\label{invall}Summary of formulas for inverse
multivectors in a left-side form of initial MV  for Clifford
algebras of dimension $p+q\le 6$. In case $p+q=6$ the second more
complicated algorithm (formula) for computation of discriminant
$s_{8+8}=\m{A}\m{G}$ in orthonormal basis is faster.}
\end{table}
The denominators, which are scalars, in these expressions are the
discriminants of respective MVs. At the same time they give the
condition for existence of the inverse MV. It was found that
discriminants exactly match expressions for determinants
calculated from matrix representation of respective  MVs.
All formulas presented in the table~\ref{invall} were checked by
explicit calculation of symbolic inverses of the most general MV
for concrete Clifford algebras with all passible signatures $(p,q)$.
For this task we wrote the {\it Mathematica} package for symbolic
calculations in the Clifford algebra~\cite{AcusDargys2017}  (to be
described elsewhere), which also can be found in {\it Mathematica}
package data base (\url{http://packagedata.net/}). The program can
simultaneously deal with a number of Clifford algebras of
different signatures $(p,q)$ what was a of great help in
computer-assisted verifications.

%%%%%%%%%%%%%%%%%%%%%%%%%%%%%%%%%%%%%%%%
\section{Outline, perspective and discussion\label{sec:5}}
From this computer-aided study we conjecture that the inverse of a
general MV of arbitrary Clifford algebra can be always  expressed
as a multilinear combination of proper number of specially
constructed self-negated products, where the initial MV may stay
in the left-most or right-most position. For the first time we
have explicitly derived the expression for inverse of general MV
for Clifford algebras of dimension $p+q=6$, which consists of
multilinear combination of just two grade-negated products. The
formulas are independent of a particular algebra signature
$(p,q)$.  They reduce to known expressions in simple cases, for
inverses of vectors, bivector or their combinations. 

The computation of inverse MV by formulas presented in
table~\ref{invall} require considerably smaller number of
multiplications because many pieces in the formulas are iterated
many times. The speed of calculation can be improved further if MV
multiplication algorithm is realized in such a way, that it can skip 
calculation of some specific grades. In this case a
huge number of operations can be avoided. For example, for
algebras of dimension $p+q=6$  the expression $\m{A}\m{A}_{{\bar
2},{\bar 3},{\bar 6}}$ consists of grades $0$, $1$, $4$ and $5$
and there is no need to calculate result for grades $2$, $3$ and $6$ since  we know
in advance that the answer will be zero.

It is easy to see that the number $m$ of multipliers in
self-negated product of discriminant in
equation~\eqref{discriminant} is determined by the dimension of
real matrices $\mathbb{R}(2)$, ${^2\mathbb{R}(2)}$,
$\mathbb{R}(4)$ etc that lie on anti-diagonals,
$n=p+q=\mathrm{const}$, of the 8-periodicity
table~\cite{Lounesto97}. For algebra with $p+q=2$ the discriminant
consists just of two multipliers, whereas for $p+q=3$ as well as
for $p+q=4$ we need to make product of four MV multipliers in
accord with dimension of a respective matrices in 8-periodicity
table. The Clifford algebras of dimensions $5$ and $6$, require
$8$ MV multivectors to construct the discriminant. Then, looking
at $8$-periodicity table one may guess that discriminant of
algebras of dimension $7$ and $8$ can possibly be constructed by
multilinear combinations from products of $16$ multivectors, etc.

The question about how many terms is required to make minimal
multilinear combination is much more difficult. We suspect that
this number might be fixed by a number of "non-eliminatible"
grades, i.~e., grades which cannot be eliminated by self-negated
product of MV with itself, which in turn would point out to the
existence of some cyclic substructures in Clifford algebra. There
are no such "non-eliminatible" grades in algebras with $p+q\le 5$.
There is only one grade, namely the grade-$4$, in algebras of
dimension $6$ and $7$, whereas the algebras with $p+q =8$ already
have two such special grades, $4$ and~$8$. It could happen that in $n=8$ case,
because of special role played by pseudoscalar (grade $8$ element
here), the both grades can be eliminated by same additive
combination consisting just of two terms. In general case,
however, multilinear combination of three terms is to be expected
for $p+q =8$ as well as for algebras of dimension~$n=9$.

From the considerations above it should be  clear that the
complexity of finding explicit expression of general MV inverse grows
rapidly with the increase of algebra dimension. For example,
direct multiplication of eight general arbitrary multivectors of
$p+q=6$ dimensional algebra requires
$(2^6)^8=281\mspace{2mu}474\mspace{3mu}976\mspace{3mu}710\mspace{4mu}656$
geometric multiplications of base elements, which fortunately can
be bypassed due to significant grade reduction in calculating the
discriminant formula as discussed in the paper. However, in
practice it does not help much when searching for correct shape of
discriminant formula. This is why reduced (shorter) MVs are always worth
to try first which, under proper choice, may greatly reduce the
number of candidates for possible combinations of multilinear
multivectors. The obtained final discriminant and inverse expressions make computations self-contained within the geometric product structure without needing to go to a matrix representation.



%\bibliographystyle{REPORT}
%\bibliography{Ghent2017}


\begin{thebibliography}{10}

\bibitem{Vaz2016}
J.~Vaz and R. {da}~Rocha,
\newblock {\em An Introduction to {C}lifford Algebras and Spinors},
\newblock Oxford University Press, Oxford, 2016.

\bibitem{Lundholm09}
D.~Lundholm and L.~Svensson,
\newblock Clifford algebra, geometric algebra, and applications,
\newblock arXiv:0907.5356v1 [math-ph], 1--117 (2009).


\bibitem{Bouma2001}
T.~A. Bouma and G.~Memowich,
\newblock Invertible homogeneous versors are blades, 2001,
\url{http://www.science.uva.nl/research/ias/ga/publications/blades.ps.gz}.

\bibitem{Semenov1991}
P.~V. Semenov,
\newblock An invertibility criterion for multivectors in a real Clifford
  algebra,
\newblock Sibirsk. Mat. Zh. {\bf 34}(4), 177--183 (1991).

\bibitem{Semenov1993A}
P.~Semenov,
\newblock On invertibility in high-dimensional {C}lifford algebras,
\newblock Found. Phys. {\bf 23}(11), 1543--1546 (1993).

\bibitem{Semenov1993B}
P.~Semenov,
\newblock On invertibility of {C}lifford algebras elements with disjoint
  supports,
\newblock in {\em Spinors, Twistors, Clifford Algebras and Quantum
  Deformations: Proceedings of the Second Max Born Symposium held near
  Wroc{\l}aw, Poland, September 1992}, edited by Z.~Oziewicz, B.~Jancewicz, and
  A.~Borowiec, pages 167--169, Springer Netherlands, Dordrecht, 1993.

\bibitem{Fletcher2002}
J.~P. Fletcher,
\newblock Clifford numbers and their inverses calculated using the matrix
  representation,
\newblock in {\em Applications of Geometric Algebra in Computer Science and
  Engineering}, edited by L.~Dorst, C.~Doran, and J.~Lasenby, pages 169--178,
  Birkh{\"a}user Boston, Boston, MA, 2002.

\bibitem{Dadbeh2011}
P.~{Dadbeh},
\newblock {Inverse and determinant in 0 to 5 dimensional {C}lifford algebra},
\newblock arXiv: 1104.0067  (Mar. 2011).

\bibitem{Shirokov2012}
D.~S. Shirokov,
\newblock Concepts of trace, determinant and inverse of {C}lifford algebra
  elements,
\newblock in {\em Proceedings of the 8th Congress of the International Society
  for Analysis, its Applications, and Computation (ISAAC)}, edited by V.~I.
  Burenkov, M.~L. Goldman, E.~B. Laneev, and V.~D. Stepanov, volume~1, pages
  187--194, Friendship University of Russia (ISBN 978-5-209-04582-3/hbk),
  Moscow, Russia, 2012,
\newblock arXiv: 1108.5447.

\bibitem{Hitzer2016}
E.~Hitzer and S.~Sangwine,
\newblock Multivector and multivector matrix inverses in real Clifford
  algebras,
\newblock Technical report, School of computer science and electronic
  engineering, Unversity of Essex, UK, Technical Report CES-534, 2016.

\bibitem{Suzuki2016}
Y.~Suzuki and N.~Yamaguchi,
\newblock Structure of inverse elements of Clifford algebra of dimension at
  most two to fifth power,
\newblock arXiv:1601.03157v1 [math-RA] , 1--14 (2016).

\bibitem{Lounesto97}
P.~Lounesto,
\newblock {\em Clifford Algebra and Spinors},
\newblock Cambridge University Press, Cambridge, 1997.

\bibitem{AcusDargys2017}
A.~Acus and A.~Dargys,
\newblock Geometric algebra Mathematica package, 2017,
 \url{https://github.com/ArturasAcus/GeometricAlgebra}.



\end{thebibliography}

\end{document}
